\begin{figure}[tb]
\centering
\begin{tikzpicture}[scale=1.05]
 \fill[derived light quantity] (0,0) rectangle (3,1.8);
 \draw[coordinate vector] (-.3,0)--(3.7,0) node[right]{$i_d$};
 \draw[coordinate vector] (0,-.3)--(0,2.4) node[above]{$i_q$};
 \draw[strong line,reference quantity,direction arrow] (0,0)--(3,0);
 \draw[strong line,reference quantity,direction arrow] (3,0)--(3,1.8);
 \draw[strong line,estimated quantity,alternative pattern,direction arrow] (0,0)--(0,1.8);
 \draw[strong line,estimated quantity,alternative pattern,direction arrow] (0,1.8)--(3,1.8);
 \node[operating point] at (0,0) {};
 \node[operating point] at (3,1.8) {};
 \node[above right] at (3,1.8) {same endpoint};
 \node[reference quantity,below] at (1.5,0) {$\mathcal C_1$: d then q};
 \node[estimated quantity,above] at (1.5,1.8) {$\mathcal C_2$: q then d};
 \node at (1.5,.9) {$\mathcal A$};
 \node[annotation,text width=66mm] at (7,.95)
  {$W'_1-W'_2=\displaystyle\iint_{\mathcal A}r_{\rm curl}\,di_d\,di_q$\\[5pt]
   Conservative field: both routes agree.\\
   Circulation: no single-valued height fits both.};
\end{tikzpicture}
\caption{Integrability as path agreement. Follow the reference quantity path and reverse
the estimated quantity path to form the counterclockwise loop used by Green's theorem.
Choosing only one measured integration path can hide inconsistency.}
\label{fig:integration}
\end{figure}
